تفاوت اصلی و بنیادین بین موجودیت (\lr{Entity}) و شیء مقداری (\lr{Value Object}) در مفهوم هویت (\lr{Identity}) نهفته است. یک \lr{Entity} دارای یک شناسه یکتا است که در طول زمان و کاملا مستقل از تغییر ویژگی‌هایش، هویت آن را ثابت نگه می‌دارد. به این معنا که اگر تمام مشخصات دو \lr{Entity} یکسان باشد اما شناسه‌های متفاوتی داشته باشند، کاملا دو شیء مجزا محسوب می‌شوند (مثل دو انسان با مشخصات ظاهری کاملا یکسان اما کدملی متفاوت). در مقابل، یک \lr{Value Object} فاقد هرگونه هویت مستقل است و تماما با ویژگی‌ها و مقادیرش تعریف می‌شود. اگر دو \lr{Value Object} دارای مقادیر برابری باشند، از نظر سیستم کاملا یکسان و قابل جایگزینی با یکدیگر هستند (مانند دو اسکناس ده هزار تومانی که برای شما تفاوتی ندارند کدام را خرج کنید). همچنین، \lr{Value Object}ها بر خلاف \lr{Entity}ها باید به صورت تغییرناپذیر (\lr{Immutable}) طراحی شوند.

در خصوص شرایط استفاده، شما زمانی باید از یک \lr{Entity} استفاده کنید که نیاز به پیگیری چرخه حیات (\lr{Lifecycle}) و تغییرات یک شیء در طول زمان دارید و تمایز آن شیء از سایر اشیاء مشابه برای منطق کسب‌وکار شما حیاتی است (مانند یک سفارش خرید یا حساب کاربری). اما زمانی باید از \lr{Value Object} استفاده کنید که هدف شما صرفا توصیف، اندازه‌گیری یا کمی‌سازی یک ویژگی در دامنه است و برایتان مهم نیست که دقیقا کدام نمونه از آن شیء را در دست دارید (مانند آدرس پستی، مبلغ پول یا یک بازه زمانی). به عنوان یک قانون کلی در رویکرد \lr{DDD}، توصیه می‌شود تا جای ممکن مفاهیم را به جای \lr{Entity} به صورت \lr{Value Object} مدل‌سازی کنید؛ زیرا به دلیل تغییرناپذیری، ایمن‌تر هستند، عوارض جانبی (\lr{Side-effects}) ندارند و مدیریت پیچیدگی سیستم را بسیار آسان‌تر می‌کنند.